% p.34 — الباب الثامن العشر في أحبّ يحبّ (مضاعف في أفعل)
\mspage{34}
\bab{الباب الثامن العشر في اَحَبَّ يُحِبُّ \quad مضاعف في اَفْعَلَ}
% Layout: six columns of continuous paradigms (see p.31 note). Col. 1 is the
% صرف صغير of أحبّ: uncontracted forms above a rule, contracted below.
% Boxed headings (title + uncontracted base form) are given as two rows.
% The lithograph writes initial hamza as a bare alif with its vowel (اَحَبَّ).
\begin{sarf}{|c|c|c|c|c|c|}
\thead{اَفْعَلَ} & \thead{ماضي معروف} & تُحِبُّ & اُحِبَّ & يُحْبَبْنَ & اَحِبُّوْا\\
اَحْبَبَ & اَحْبَبَ & تُحِبَّانِ & اُحِبَّا & تُحَبُّ & اَحِبِّيْ\\
يُحْبِبُ & اَحَبَّ & يُحْبِبْنَ & اُحِبُّوْا & تُحَبَّانِ & اَحِبَّا\\
اِحْبَابًا & اَحَبَّا & تُحِبُّ & اُحِبَّتْ & تُحَبُّوْنَ & اَحْبِبْنَ\\
مُحْبِبٌ & اَحَبُّوْا & تُحِبَّانِ & اُحِبَّتَا & تُحَبِّيْنَ & \thead{نهي حاضر}\\
اُحْبِبَ & اَحَبَّتْ & تُحِبُّوْنَ & اُحْبِبْنَ & تُحَبَّانِ & لَا تُحْبِبْ\\
يُحْبَبُ & اَحَبَّتَا & تُحِبِّيْنَ & اُحْبِبْتَ & تُحْبَبْنَ & لَا تُحِبَّ\\
اِحْبَابًا & اَحْبَبْنَ & تُحِبَّانِ & اُحْبِبْتُمَا & اُحَبُّ & لَا تُحِبَّا\\
مُحْبَبٌ & اَحْبَبْتَ & تُحْبِبْنَ & اُحْبِبْتُمْ & نُحَبُّ & لَا تُحِبُّوْا\\
اَحْبِبْ & اَحْبَبْتُمَا & اُحِبُّ & اُحْبِبْتِ & \thead{اسم مفعول} & لَا تُحِبِّيْ\\
لَا تُحْبِبْ & اَحْبَبْتُمْ & نُحِبُّ & اُحْبِبْتُمَا & مُحْبَبٌ & لَا تُحِبَّا\\
اَحَبَّ & اَحْبَبْتِ & \thead{اسم فاعل} & اُحْبِبْتُنَّ & مُحَبٌّ & لَا تُحْبِبْنَ\\
يُحِبُّ & اَحْبَبْتُمَا & مُحْبِبٌ & اُحْبِبْتُ & مُحَبَّانِ & \thead{مضارع بلم}\\
اِحْبَابًا & اَحْبَبْتُنَّ & مُحِبٌّ & اُحْبِبْنَا & مُحَبُّوْنَ & لَمْ يُحْبِبْ\\
مُحِبٌّ & اَحْبَبْتُ & مُحِبَّانِ & \thead{مضارع مجهول} & مُحَبَّةٌ & لَمْ يُحِبَّ\\
اُحِبَّ & اَحْبَبْنَا & مُحِبُّوْنَ & يُحْبَبُ & مُحَبَّتَانِ & لَمْ يُحِبَّا\\
يُحَبُّ & \thead{مضارع معروف} & مُحِبَّةٌ & يُحَبُّ & مُحَبَّاتٌ & لَمْ يُحِبُّوْا\\
اِحْبَابًا & يُحْبِبُ & مُحِبَّتَانِ & يُحَبَّانِ & \thead{أمر حاضر} & لَمْ تُحِبَّ\\
مُحَبٌّ & يُحِبُّ & مُحِبَّاتٌ & يُحَبُّوْنَ & اَحْبِبْ & لَمْ تُحِبَّا\\
اَحِبَّ & يُحِبَّانِ & \thead{ماضي مجهول} & تُحَبُّ & اَحِبَّ & لَمْ يُحْبِبْنَ\\
لَا تُحِبَّ & يُحِبُّوْنَ & اُحْبِبَ & تُحَبَّانِ & اَحِبَّا & لَمْ تُحِبَّ\\
\end{sarf}
